\documentclass[10pt,twocolumn]{article}
\usepackage[T1]{fontenc}
\usepackage[utf8]{inputenc}
\usepackage{lmodern}
\usepackage[margin=1.8cm,top=2.4cm,bottom=2cm,columnsep=0.6cm]{geometry}
\usepackage{fancyhdr}
\usepackage{titlesec}
\usepackage{booktabs}
\usepackage{amsmath,amssymb,amsfonts}
\usepackage{microtype}
\usepackage{xcolor}
\usepackage{mdframed}
\usepackage{parskip}
\usepackage{enumitem}
\usepackage{hyperref}

\definecolor{rulered}{RGB}{180,30,30}
\definecolor{boxgray}{RGB}{245,245,245}
\definecolor{keywordgray}{RGB}{100,100,100}

\pagestyle{fancy}
\fancyhf{}
\fancyhead[L]{\textsc{\small Claude Mag}}
\fancyhead[C]{\small\textit{Machine Learning Research Digest}}
\fancyhead[R]{\small 2026-09-23}
\fancyfoot[C]{\small\thepage}
\renewcommand{\headrulewidth}{0.8pt}

\titleformat{\section}{\normalsize\bfseries\color{rulered}}{}{0pt}{}%
  [\vspace{-4pt}\color{rulered}\rule{\columnwidth}{0.4pt}]
\titlespacing{\section}{0pt}{8pt}{4pt}

\hypersetup{colorlinks=true,urlcolor=rulered,linkcolor=black}

\begin{document}
\setlength{\parindent}{0pt}
\setlength{\parskip}{4pt}

{\small\color{keywordgray}\textbf{Keywords:}
  humanoid locomotion \textbullet{}
  constrained reinforcement learning \textbullet{}
  smoothness \textbullet{}
  sim-to-real}\par\vspace{4pt}

{\LARGE\bfseries\raggedright
  Smoothness as a Constraint for Stable
  Humanoid Locomotion}\par\vspace{4pt}

{\large\itshape\raggedright
  DeCap Decouples Upper- and Lower-Body Smoothness into Separate
  Constraint Groups, Cutting Real-Robot Action Rate by 2.50$\times$
  Without Reward Retuning Across Terrains}\par\vspace{6pt}

{\small\textbf{By} Utsav Panchal, Denis Kleyko, Unal Artan \&
  Amy Loutfi}\par\vspace{2pt}

{\small\color{keywordgray}arXiv:2609.24552 \textbullet{} September 2026}
\par\vspace{2pt}\noindent{\color{rulered}\rule{\columnwidth}{0.6pt}}\vspace{6pt}

Reward-based smoothness penalties compete with locomotion objectives and
treat the body uniformly. DeCap reformulates whole-body smoothness as
explicit upper- and lower-body constraints with a bounded barrier penalty,
reducing upper-body action rate by 2.50$\times$ and acceleration by
2.18$\times$ on real humanoids while transferring to new terrains without
reward retuning.

\begin{mdframed}[backgroundcolor=boxgray,linecolor=rulered,linewidth=1pt,
    innerleftmargin=6pt,innerrightmargin=6pt,
    innertopmargin=6pt,innerbottommargin=6pt]
\textbf{\small AT A GLANCE}\par\vspace{4pt}
\begin{description}[leftmargin=0pt,labelsep=4pt,itemsep=2pt,parsep=0pt,
    font=\small\bfseries]
  \item[Problem:] Reward-based smoothness competes with task rewards
    and treats upper and lower body uniformly, ignoring their different
    physical roles in humanoid stability.
  \item[Why it matters:] Jerky upper-body motions destabilize center-of-mass
    trajectories; over-smoothing the lower body prevents reactive stepping
    on uneven terrain.
  \item[Method:] DeCap: separate tight/loose bounded barrier constraints
    for upper/lower body, formulated as an RL constraint term.
  \item[Result:] 2.50$\times$ upper-body action-rate reduction,
    2.18$\times$ acceleration reduction; terrain transfer without reward
    retuning.
  \item[Evidence:] Real humanoid whole-body control evaluation on flat
    ground, stairs, slopes, and stepping stones; 4-axis ablation study.
\end{description}
\end{mdframed}

\section{The Big Picture}

Humanoid robots are expected to walk, carry objects, and navigate
complex terrains while remaining stable and human-like in motion.
Reinforcement learning has enabled impressive locomotion policies,
but the physical quality of the motion---smoothness, absence of
oscillations, predictability---has been an afterthought, typically
handled by adding penalty terms to the reward function.

This paper argues that treating smoothness as a reward penalty is
the wrong abstraction for two reasons:

\textbf{First, it creates objective competition.} A smoothness penalty
competes with speed rewards and stability bonuses during gradient descent.
The resulting trade-off must be re-tuned for every new terrain, speed
target, and robot variant---a laborious process that does not transfer.

\textbf{Second, it treats the body uniformly.} The upper body (torso,
arms) must be tightly regulated to prevent tipping---inertial forces
from fast arm motion directly perturb the center-of-mass. The lower body
(hips, knees, ankles) must remain reactive to ground-contact perturbations;
over-smoothing the legs leads to falls on uneven terrain.

DeCap (Decoupled Constraint-aware Policy) addresses both problems by
reformulating smoothness as explicit, body-part-specific constraints
within the RL objective.

\section{Why Existing Approaches Fall Short}

\textbf{Reward-based action-rate penalties} are soft: the optimizer
can trade smoothness for task performance whenever doing so increases
total reward. The penalty coefficient requires extensive per-terrain tuning.

\textbf{Lipschitz-constrained policies} bound the policy's output
sensitivity globally but cannot differentiate between upper and lower
body requirements.

\textbf{Model predictive control (MPC)} provides smooth trajectories by
design but requires accurate dynamics models and fails on highly uncertain
terrain contacts.

\textbf{Trajectory imitation} from motion capture data inherits human
smoothness but is brittle to out-of-distribution perturbations and
terrain types.

\section{Core Mechanism}

\textbf{Decoupled Constraint Groups.} DeCap partitions joints into
upper-body set $\mathcal{U}$ (torso, arms, neck) and lower-body set
$\mathcal{L}$ (hips, knees, ankles), reflecting their distinct physical
roles. Smoothness is measured as joint action rate:
\[
\dot{a}_k^{(j)} = \frac{a_k^{(j)} - a_{k-1}^{(j)}}{\Delta t}
\]

Two separate constraint limits are set: tight $\kappa_U$ for the upper
body (derived from human motion capture statistics), loose $\kappa_L$
for the lower body (set to allow reactive stepping).

\textbf{Bounded Barrier Penalty.} Instead of a standard log-barrier
(which diverges at the constraint limit and causes training instability),
DeCap uses a bounded barrier that activates proactively in the range
$[\alpha\kappa,\kappa)$ and saturates at the limit:
\[
\phi(c,\kappa) = \begin{cases}
0 & c \le \alpha\kappa \\[3pt]
\dfrac{(c-\alpha\kappa)^2}{\kappa-\alpha\kappa} & \alpha\kappa < c < \kappa \\[3pt]
\kappa - \alpha\kappa & c \ge \kappa
\end{cases}
\]
where $c = \|\dot{\mathbf{a}}_k\|_2$ is the action-rate norm and
$\alpha\in(0,1)$ is the activation threshold fraction.

\section{Technical Formulation}

Upper-body constraint barrier reward term:
\[
R_{\mathrm{smooth}}^U
= -\lambda_U\cdot\phi\!\left(\|\dot{\mathbf{a}}_k^U\|_2,\;\kappa_U\right)
\]

Lower-body constraint barrier:
\[
R_{\mathrm{smooth}}^L
= -\lambda_L\cdot\phi\!\left(\|\dot{\mathbf{a}}_k^L\|_2,\;\kappa_L\right)
\]

Second-order (acceleration) constraint for the upper body:
\[
\ddot{a}_k^{(j)} = \frac{\dot{a}_k^{(j)} - \dot{a}_{k-1}^{(j)}}{\Delta t},
\quad
R_{\mathrm{acc}}^U
= -\lambda_{\mathrm{acc}}\cdot\phi\!\left(
  \|\ddot{\mathbf{a}}_k^U\|_2,\;\kappa_U^{\mathrm{acc}}\right)
\]

Full RL objective:
\[
\mathcal{L}(\theta) = \mathcal{L}_{\mathrm{task}}(\theta)
  + R_{\mathrm{smooth}}^U + R_{\mathrm{smooth}}^L + R_{\mathrm{acc}}^U
\]

\section{Learning or Inference Procedure}

\begin{enumerate}[itemsep=2pt,parsep=0pt]
  \item Partition robot URDF joints into $\mathcal{U}$ and $\mathcal{L}$.
  \item Set $\kappa_U$ from human motion-capture action-rate statistics;
    set $\kappa_L$ empirically to allow reactive stepping.
  \item Train PPO in Isaac Gym on flat ground with domain randomisation.
  \item Deploy on real humanoid; evaluate on flat ground and 3 novel
    terrains (stairs, slopes, stepping stones) without reward retuning.
\end{enumerate}

\section{What the Guarantee Says}

The bounded barrier provides a formal advantage over log-barriers: the
gradient is bounded everywhere (including at the constraint limit),
preventing catastrophic gradient magnitudes that destabilise RL training.
For $c \ge \kappa$, the barrier saturates at $\kappa - \alpha\kappa$, a
finite value that continues to penalise violations while remaining numerically
stable. The proactive activation at $\alpha\kappa$ teaches the policy to
stay away from the limit before it is reached.

\section{Experimental Findings}

\begin{itemize}[itemsep=2pt,parsep=0pt]
  \item \textbf{Upper-body action rate:} 2.50$\times$ reduction over
    reward-based smoothness baseline on flat ground
  \item \textbf{Upper-body acceleration:} 2.18$\times$ reduction
  \item \textbf{Lower-body smoothness:} Marginal improvement;
    reactive range preserved
  \item \textbf{Transient motion:} Reduced oscillations in torso and
    arm joints during gait transitions
  \item \textbf{Terrain transfer:} Fixed constraints transfer to stairs,
    slopes, and stepping stones without reward retuning; reward-based
    baselines require significant coefficient adjustment on each terrain
\end{itemize}

\section{Ablations and Interpretation}

\textbf{No decoupling (single constraint group):} Upper-body smoothness
improves, but lower-body reactivity degrades on uneven terrain, leading
to falls; confirms that separate limits are necessary.

\textbf{Standard log-barrier:} Training instability due to gradient
explosions near the constraint limit; DeCap's bounded barrier resolves this.

\textbf{Reward term instead of constraint:} Smoothness improvement is
$\approx$40\% of DeCap's at matched task-reward coefficient, confirming
that the constraint formulation provides direct, stable control over the
smoothness quantities.

\textbf{$\alpha$ sensitivity:} $\alpha=0.7$ provides the best balance;
smaller $\alpha$ (earlier activation) over-constrains and reduces
task reward by $\approx 5\%$.

\section*{Reference}

Utsav Panchal, Denis Kleyko, Unal Artan, Amy Loutfi.
\textbf{Smoothness as a Constraint for Stable Humanoid Locomotion}.
arXiv:2609.24552, September 2026.\\
\url{https://arxiv.org/abs/2609.24552}

\end{document}
